\section{泛函的定义}
\label{sec:lfunc-definition}

函数空间中的元素是整个函数，我们对它所做的许多操作最后只得到一个数。把这类操作统一起来，就得到泛函的概念。本节先看几个常见的例子，逐一核对它们是否线性、在哪个空间上有定义；再研究线性泛函什么时候连续，并用一个数度量它把输入误差放大多少。全章记标量域为$\mathbb F$，它是$\mathbb R$或$\mathbb C$。

\begin{definition}[泛函]
\label{def:lfunc-functional}
设$X$是$\mathbb F$上的线性空间。映射$\ell:X\to\mathbb F$称为$X$上的\term{泛函}（functional）。若对一切$x,y\in X$和$\alpha,\beta\in\mathbb F$都有$\ell(\alpha x+\beta y)=\alpha\ell(x)+\beta\ell(y)$，称$\ell$为\term{线性泛函}（linear functional）。
\end{definition}

有限维空间上的线性泛函由它在一组基上的取值决定，可以写成坐标的线性组合，这在\BookChapterRef{chap:linear-representation}{卷二第2章“线性空间、映射与坐标”}中已经讨论过。函数空间是无限维的，新问题在于：线性泛函未必连续，读数是否稳定要看空间所配的范数。

\subsection{常见泛函}
\label{sec:lfunc-examples}

下面四个例子分别代表加权平均、坐标读取、点上取值和能量。每个例子都要先说清它作用在哪个空间上，因为同一条规则在一个空间上有意义，在另一个空间上可能根本没有定义。

\paragraph{加权积分}
设$(\Omega,\mathcal F,\mu)$是测度空间，$1\leq p\leq\infty$，$q$是\term{共轭指数}（conjugate exponent），即$1/p+1/q=1$，约定$p=1$时$q=\infty$、$p=\infty$时$q=1$。对固定的$g\in L^q(\mu)$，令
\begin{equation}
L_g(f)=\int_\Omega fg\,d\mu\qquad(f\in L^p(\mu)).
\label{eq:lfunc-weighted-integral}
\end{equation}
由第~\ref{chap:normed-spaces}章的Hölder不等式，$\int\lvert fg\rvert\,d\mu\leq\lVert f\rVert_p\lVert g\rVert_q<\infty$，积分有定义。线性来自积分的线性：$L_g(\alpha f_1+\beta f_2)=\alpha L_g(f_1)+\beta L_g(f_2)$。

$L^p$的元素是几乎处处相等的函数组成的等价类，所以还要检查读数不依赖代表的选取。若$f=\tilde f$几乎处处成立，则$fg=\tilde fg$几乎处处成立，两者的积分相同。积分对零测集不敏感，这正是加权积分能在等价类上使用的原因。

\paragraph{正交系数}
设$H$是Hilbert空间，$\{e_k\}$是标准正交系。第$k$个系数$f\mapsto\langle e_k,f\rangle$是$H$上的泛函；按第~\ref{chap:inner-product-spaces}章的约定，内积对第二个变量线性，所以它是线性泛函。在$L^2(0,1)$中，它就是以$\overline{e_k}$为权重的加权积分。

\paragraph{点评估}
设$K$是紧度量空间，$C(K)$是$K$上的连续函数空间。固定$x_0\in K$，令$\delta_{x_0}(f)=f(x_0)$，称为$x_0$处的\term{点评估泛函}（point evaluation functional）。它显然是线性的。

点评估要求逐点取值，只在由真正的函数组成的空间上有定义。在Lebesgue测度下的$L^p(0,1)$中，单点$\{x_0\}$的测度为零，$f$与$f+\mathbf{1}_{\{x_0\}}$属于同一个等价类，在$x_0$处的值却不同。所以“读出$x_0$处的值”不是$L^p$上的映射，谈不上线性或连续。积分与点评估的区别就在这里：前者对零测集不敏感，后者恰恰只看一个零测集。

\paragraph{平方范数}
在赋范空间上，$f\mapsto\lVert f\rVert^2$也把函数变成一个数，例如$L^2$中的能量$\int\lvert f\rvert^2$。它不是线性的：$\lVert2f\rVert^2=4\lVert f\rVert^2\neq2\lVert f\rVert^2$（$f\neq0$时）。非线性泛函在优化中很常见，如以函数为未知量的能量与损失；它们的变分和极小化可参考\BookChapterRef{chap:functional-optimization-variation}{卷五第11章“泛函优化与变分”}。本章只研究线性泛函。

\subsection{泛函范数}
\label{sec:lfunc-norm}

读数是否稳定，就是输入的小误差是否只引起读数的小误差。对线性泛函，这个要求有简单的等价形式。设$X$是赋范空间。由第~\ref{chap:metric-space-maps}章的距离判据，$\ell$在$x_0$处连续，意思是对每个$\varepsilon>0$存在$\delta>0$，使$\lVert x-x_0\rVert<\delta$时$\lvert\ell(x)-\ell(x_0)\rvert<\varepsilon$。

\begin{theorem}[线性泛函的连续性]
\label{thm:lfunc-continuity-bounded}
设$X$是赋范空间，$\ell$是$X$上的线性泛函。以下三条等价：
\begin{enumerate}
\item $\ell$在$0$处连续；
\item $\ell$在$X$上处处连续；
\item 存在常数$C\geq0$，使$\lvert\ell(x)\rvert\leq C\lVert x\rVert$对一切$x\in X$成立。
\end{enumerate}
\end{theorem}

\begin{proof}
由3推2：$\lvert\ell(x)-\ell(x_0)\rvert=\lvert\ell(x-x_0)\rvert\leq C\lVert x-x_0\rVert$，取$\delta=\varepsilon/(C+1)$即可。由2推1是显然的。

由1推3：取$\varepsilon=1$，由$\ell(0)=0$，存在$\delta>0$使$\lVert z\rVert<\delta$时$\lvert\ell(z)\rvert<1$。对$x\neq0$，令$z=\frac{\delta}{2\lVert x\rVert}x$，则$\lVert z\rVert=\delta/2<\delta$，于是由线性
\[
\lvert\ell(x)\rvert=\frac{2\lVert x\rVert}{\delta}\lvert\ell(z)\rvert<\frac{2}{\delta}\lVert x\rVert .
\]
$x=0$时不等式两边都是零。取$C=2/\delta$。
\end{proof}

线性在证明中起了两次作用：把任意点处的误差平移成原点附近的误差，把原点附近的控制按比例放大到所有向量。满足第3条的线性泛函称为\term{有界线性泛函}（bounded linear functional）；由定理~\ref{thm:lfunc-continuity-bounded}，对线性泛函而言“有界”与“连续”是同一回事。这里的有界指在单位球上有界，而不是在整个空间上有界，后者只有零泛函才满足。

\begin{definition}[泛函范数]
\label{def:lfunc-functional-norm}
设$\ell$是赋范空间$X$上的有界线性泛函，令
\[
\lVert\ell\rVert=\sup_{\lVert x\rVert\leq1}\lvert\ell(x)\rvert,
\]
称为$\ell$的\term{泛函范数}（norm of a functional）。
\end{definition}

由线性，$\lvert\ell(x)\rvert\leq\lVert\ell\rVert\lVert x\rVert$对一切$x$成立，并且$\lVert\ell\rVert$是使这一不等式成立的最小常数$C$。它衡量的是读数对输入误差的最大放大倍数。

\paragraph{点评估的范数依赖空间}
在$C([0,1])$上配一致范数$\lVert f\rVert_\infty=\max\lvert f\rvert$，由$\lvert f(0)\rvert\leq\lVert f\rVert_\infty$得$\lVert\delta_0\rVert\leq1$，取$f\equiv1$得$\lVert\delta_0\rVert=1$。点评估在一致范数下是稳定的读数。

同一个集合$C([0,1])$改配积分范数，结论完全不同。取贴着左端点的尖峰
\[
\tau_n(t)=\max(1-nt,0)\qquad(t\in[0,1]),
\]
它连续，$\tau_n(0)=1$，而直接积分得
\[
\lVert\tau_n\rVert_1=\int_0^{1/n}(1-nt)\,dt=\frac{1}{2n},\qquad
\lVert\tau_n\rVert_2^2=\int_0^{1/n}(1-nt)^2\,dt=\frac{1}{3n}.
\]
两种积分范数都趋于零，读数却始终是$1$，所以不存在常数$C$使$\lvert f(0)\rvert\leq C\lVert f\rVert_2$对一切连续$f$成立。在$(C([0,1]),\lVert\cdot\rVert_2)$上，$\delta_0$是线性的、处处有定义，却不连续。

\paragraph{加权积分的范数也依赖空间}
在$L^2(0,1)$上取$g=\mathbf{1}_{[0,1/2]}$，$L_g(f)=\int_0^{1/2}f$。Hölder不等式给出$\lVert L_g\rVert\leq\lVert g\rVert_2=1/\sqrt2$；取$f=g$，$L_g(g)=1/2$而$\lVert g\rVert_2=1/\sqrt2$，比值为$1/\sqrt2$，所以$\lVert L_g\rVert=1/\sqrt2$。同一规则在$C([0,1])$的一致范数下，由$\lvert\int_0^{1/2}f\rvert\leq\lVert f\rVert_\infty/2$及$f\equiv1$得范数为$1/2$。读数规则相同，范数不同，因为衡量输入误差的方式不同。是否稳定、放大多少，都依赖所选的空间与范数。
